%% kapa: exam=1 round=32 year=2021 period=1 form=A subject=civil_law subject_ko=민법 seq=40 number=40
%% source: Q-Net 감정평가사 기출문제 (cst00309) artlSeq=5208397 file=extracted/1008/5208397/A형/2021년도 제32회 감정평가사 1차 1교시 A형.pdf
%% conversion: kordoc 4.24.51 -> tools/gen1.py (auto); review=proofread-sonnet; flags=none
\begin{question}{40}
전세권을 목적으로 하는 저당권에 관한 설명으로 옳지 않은 것은? (다툼이 있으면 판례에 따름)
\begin{choices}
  \item 저당권설정자는 저당권자의 동의 없이 전세권을 소멸하게 하는 행위를 하지 못한다.
  \item 전세권의 존속기간이 만료된 경우 저당권자는 전세권 자체에 대해 저당권을 실행할 수 있다.
  \item 전세권의 존속기간이 만료되면 저당권자는 전세금반환채권에 대하여 물상대위할 수 있다.
  \item 전세금반환채권은 저당권의 목적물이 아니다.
  \item 전세권이 기간만료로 소멸한 경우 전세권설정자는 원칙적으로 전세권자에 대하여만 전세금반환의무를 부담한다.
\end{choices}
\end{question}
\answer{2}
